% Esquisse d'un programme (cote 119, copy 2): archive cote 119, site page 122 (PDF page 123), typescript page 44.
% typescript. Draft compared: esquisse/pages/esq_119_046.tex (copy 1, site 46).
% AI-transcribed by Claude Opus 5.5 (claude-opus-5-5, adaptive thinking, effort high), Message Batches API, batch esq119_full (msgbatch_015xqDHkTumDysRF1qngZhRJ), 2026-11-13; not yet checked by eye.
\scanpage{119/122}
\mspage{44}

différente de celle qui me tenait et tient à coeur. Il introduit
certes des catégories dérivées dans divers contextes non commu-%
tatifs, mais sans entrer dans la question des structures internes
essentielles d'une telle catégorie, laissée ouverte également
dans le cas commutatif par Verdier, et après lui par Illusie. De
même, la question de mettre le doigt sur les ``coefficients''
naturels pour un formalisme cohomologique non commutatif, au delà
des champs (qu'on devrait appeler 1-champs) étudiés dans le livre
de Giraud, restait ouverte - ou plutôt, les intuitions riches et
précises qui y répondent, puisées dans des exemples nombreux pro-%
venant de la géométrie algébrique notamment, attendent toujours
un langage précis et souple pour leur donner forme.

Je reviens sur certains aspects de ces questions de
fondements en 1975, à l'occasion (je crois me souvenir) d'une
correspondance avec Larry Breen (\deleted{deux}\inserted{trois} lettres de cette corres-%
pondance seront reproduites en appendice au Chap. I du volume 1,
``Histoire de Modèles'', de la Poursuite des Champs). A ce moment
apparaît l'intuition que les $\infty$-groupoïdes doivent constituer des
modèles, particulièrement adéquats, pour les types d'homotopie,
les n-groupoïdes correspondant aux types d'homotopie \uline{tronqués}
(avec $\pi_i = 0$ pour $i > n$). Cette même intuition, par des voies
très différentes, a été retrouvée par Ronnie Brown à Bangor et
certains de ses élèves, mais en utilisant une notion de $\infty$-groupoïde
assez restrictive (qui, parmi les types d'homotopie 1-connexes,
ne modélise que les produits d'espacesd'Eilenberg-Maclane). C'est
stimulé par une correspondance à bâtons rompus avec Ronnie Brown,
que j'ai finalement repris une réflexion, commençant par un essai
de définition d'une notion de $\infty$-groupoïde plus large (rebaptisé
par la suite ``champ en $\infty$-groupoïdes'' ou simplement ``champ'', sous-%
entendu : sur le topos ponctuel), et qui de fil en aiguille m'a
amené à la Poursuite des Champs. Le volume ``Histoire de Modèles''
y constitue d'ailleurs une digression entièrement imprévue par
rapport au propos initial (les fameux champs étant provisoirement
oubliés, et n'étant prévus réapparaître que vers les pages 1 000
environ ...).

\margin{122}
